% gz-p10-07-1: 构造法证 ln(1+x) < x（x > 0）——y=ln(1+x) 和 y=x 的对比
\documentclass[tikz,border=4pt]{standalone}
\usepackage{ctex}
\usepackage{amsmath}
\usepackage{pgfplots}
\pgfplotsset{compat=1.18}

\begin{document}
\begin{tikzpicture}
\begin{axis}[
  width=10cm, height=7cm,
  axis lines=center,
  xlabel={$x$}, ylabel={$y$},
  every axis x label/.style={at={(axis cs:3.6,0)}, anchor=west},
  every axis y label/.style={at={(axis cs:0,3.7)}, anchor=south},
  xmin=-1.0, xmax=3.6,
  ymin=-1.5, ymax=3.7,
  xtick={1,2,3},
  ytick={1,2,3},
  tick style={color=black},
  ticklabel style={font=\small},
  clip=true,
]
  % 蓝线：y = ln(1+x)（x > -1）
  \addplot[blue, thick, domain=-0.95:3.5, samples=120] {ln(1+x)};
  \node[right, blue] at (axis cs:2.5, 1.0) {$y=\ln(1+x)$};

  % 红线：y = x
  \addplot[red, thick, domain=-1.0:3.5] {x};
  \node[right, red] at (axis cs:2.0, 2.4) {$y=x$};

  % 交点：x = 0
  \addplot[only marks, mark=*, mark size=2.5pt, black] coordinates {(0,0)};

  % 注释：x > 0 时红在蓝上方
  \node[font=\footnotesize, align=left, red] at (axis cs:1.5, -0.8)
    {当 $x>0$，红 $>$ 蓝\\即 $\ln(1+x)<x$};

  % 阴影示意 x>0 区域
  \draw[red, dashed, thick] (axis cs:1.5, 1.5) -- (axis cs:1.5, 0.916);
  \node[right, font=\footnotesize] at (axis cs:1.5, 1.2) {$x-\ln(1+x)>0$};
\end{axis}
\end{tikzpicture}
\end{document}
